\subsection{特征标与共轭类}

令 \(\varrho\) 为 G 在有限维希尔伯特空间 E 中的酉表示。\(\varrho\) 的特征标满足 \(|\chi_{\varrho}| \leqslant \dim(E)\) 。令 \((e_i)_{i \in I}\) 为 E 的标准正交基；特征标 \(\chi_{\varrho}\) 是对角矩阵系数 \(g \mapsto \langle e_i | \varrho(g)e_i \rangle\)（属于 \(\varrho\)）之和。特别地，\(\varrho\) 的特征标是 G-有限函数；若 \(\varrho\) 不可约，则 \(\chi_{\varrho} \in \varrho_{\mathrm{G}}(\varrho)\) 。

有如下公式

\[
\chi_ {\varrho_ {1} \oplus \varrho_ {2}} = \chi_ {\varrho_ {1}} + \chi_ {\varrho_ {2}}, \quad \chi_ {\varrho_ {1} \otimes \varrho_ {2}} = \chi_ {\varrho_ {1}} \chi_ {\varrho_ {2}}, \quad \chi_ {\breve {\varrho}} = \chi_ {\overline {{\varrho}}} = \overline {{\chi}} _ {\varrho}
\]

（参见 A, VIII, 第 388--389 页）。

命题 8. --- 有

\begin{enumerate}
\def\labelenumi{(\arabic{enumi})}
\tightlist
\item
  \(\chi_{\pi}*\chi_{\sigma} = 0\) 对于 \(\pi, \sigma\) 均属于 \(\widehat{\mathrm{G}}\) 且 \(\pi \neq \sigma\) 的情形成立，
\end{enumerate}

\[
\chi_ {\pi} * \chi_ {\pi} = \frac {1}{\dim (\pi)} \chi_ {\pi} \quad\text{对于每个 }\pi\in\widehat {\mathrm{G}}.\tag{2}
\]

令 \(\pi\) 和 \(\sigma\) 为 G 的不可约表示。有

\[
\dim (\pi) \left(\chi_ {\pi} * \chi_ {\sigma}\right) = \dim (\pi) \boldsymbol {\gamma} _ {G} \left(\chi_ {\pi}\right) \chi_ {\sigma}
\]

（V, 第 407 页引理 4）。函数 \(\dim(\pi)\gamma_{\mathrm{G}}(\chi_{\pi})\chi_{\sigma}\) 是将 \(\chi_{\sigma}\) 投影到 \(\overline{\pi}\)-同型分量 \(\gamma_{\mathrm{G}}\) 上的正交投影（V, 第 465 页推论 2），即投影到 \(\varrho_{\mathrm{G}}(\pi)\) 上（V, 第 463 页推论 4）。由于 \(\chi_{\sigma}\) 属于 \(\varrho_{\mathrm{G}}(\sigma)\)，结论由 V, 第 462 页定理 1 得出。

引理 2. --- 等价关系“\(x \in G\) 且 \(y \in G\) 且 \(x\) 与 \(y\) 共轭”在 G 中的图像是闭集。

事实上，该图像是从紧空间 G × G 到自身、由 \((x, y) \mapsto (x, yxy^{-1})\) 定义的连续映射的像。

记 G \(^{\sharp}\) 为赋予商拓扑的 G 的共轭类空间；由引理 2 以及 TG, I, 第 78 页命题 8，它是紧空间。令 \(\varpi\) : G → G \(^{\sharp}\) 为典范投影。由从 \(\mathcal{C}(G^{\sharp})\) 到 \(\mathcal{C}(G)\) 的映射 \(f \mapsto f \circ \varpi\)，将星代数 \(\mathcal{C}(G^{\sharp})\) 识别为 \(\mathcal{C}(G)\) 中由连续中心函数组成的星子代数。

G \(^{\sharp}\) 上的测度与 G 上的中心测度相对应（V, 第 402 页定义 1）。

于是，在 \(\mathcal{C}(\mathrm{G}^{\sharp})\) 上由 \(f\mapsto\int_{\mathrm{G}}f\) 定义的线性泛函是 \(\mathrm{G}^{\sharp}\) 上质量为 1 的正测度，记为 \(\mu^{\sharp}\)。对于每个 \(p\in[1,+\infty]\)，记 \(\mathcal{L}^{p}(\mathrm{G}^{\sharp})\)（分别记 \(\mathrm{L}^{p}(\mathrm{G}^{\sharp})\)）为 \(\mathcal{L}_{\mathbf{C}}^{p}(\mathrm{G}^{\sharp},\mu^{\sharp})\)（分别为 \(\mathrm{L}_{\mathbf{C}}^{p}(\mathrm{G}^{\sharp},\mu^{\sharp})\)）。我们将 \(\mathrm{L}^{p}(\mathrm{G}^{\sharp})\) 识别为 \(\mathcal{C}(\mathrm{G}^{\sharp})\) 在 \(\mathrm{L}^{p}(\mathrm{G})\) 中的闭包；特别地，它是 \(\mathrm{L}^{p}(\mathrm{G})\) 的闭子空间。

我们还记 \(\Theta(G^{\sharp}) = \Theta(G) \cap \mathcal{C}(G^{\sharp})\) 为 G 的中心矩阵系数空间。它是 \(\mathcal{C}(G^{\sharp})\) 的含单位对合子代数。

当 G 为李群时，在 LIE, IX, 第 71 页中，空间 \(\Theta(G^{\sharp})\) 也记为 \(Z\Theta(G)\)。

引理 3. --- 令 \(\pi\) 为 G 的不可约酉表示。向量空间 \(\mathrm{L}^2 (\mathrm{G}^\sharp)\cap \varrho_{\mathrm{G}}(\pi)\) 是一维的，并由 \(\pi\) 的特征标生成。

考虑 G 的酉表示 \(\sigma\)，它作用在空间 \(\varrho_{\mathrm{G}}(\pi)\) 上，并由 \(\sigma(g) = \varrho_{\mathrm{G}}(g, g)\) 定义。由于 G × G 的酉表示 \(\varrho_{\mathrm{G}}(\pi)\) 同构于 \(\overline{\pi} \boxtimes \pi\)（V, 第 462 页命题 5），其特征标为

\[
\chi_ {\sigma} (g) = \chi_ {\overline {{\pi}} \boxtimes \pi} (g, g) = | \chi_ {\pi} (g) | ^ {2}
\]

对所有 \(g \in G\) 都成立。\(\mathrm{L}^2 (\mathrm{G}^\sharp) \cap \varrho_{\mathrm{G}}(\pi)\) 是该表示的不变元素子空间，且其维数等于

\[
\int_ {\mathrm{G}} \chi_ {\sigma} (g) d \mu (g) = \int_ {\mathrm{G}} | \chi_ {\pi} (g) | ^ {2} d \mu (g) = 1
\]

（V, 第 466 页推论 3 以及第 459 页推论 3）。由于 \(\pi\) 的特征标是 \(L^2(G^\sharp) \cap \varrho_G(\pi)\) 的非零元素，它就是该空间的一个基。

命题 9. --- G 的不可约酉表示的特征标族 \((\chi_{\pi})_{\pi \in \widehat{\mathbf{G}}}\) 是 \(\mathrm{L}^2 (\mathrm{G}^\sharp)\) 的标准正交基。

根据彼得–外尔定理（V, 第 462 页定理 1），不变元素构成闭子空间 \(L^{2}(G^{\sharp})\)，它是 \(L^{2}(G)\) 的子空间，由 G 的酉表示 \(g \mapsto \varrho_{\mathrm{G}}(g, g)\) 的不变元素构成；它是各子空间 \(L^{2}(G^{\sharp}) \cap \varrho_{\mathrm{G}}(\pi)\)（\(\pi \in \widehat{G}\)）的希尔伯特直和。因此命题由前一引理和 V, 第 459 页推论 3 得出。

推论 1. --- 向量空间 \(\Theta(\mathrm{G}^{\sharp})\) 在 \(\mathcal{C}(\mathrm{G}^{\sharp})\) 和 \(\mathrm{L}^{2}(\mathrm{G}^{\sharp})\) 中都是稠密的。

第二个断言由命题 9 得出。对于第一个断言，考虑 G 的线性表示 \(\varrho\)，它作用在巴拿赫空间 \(\mathcal{C}(\mathrm{G})\) 上，并由下式定义：

\[
\varrho (g) f (x) = f (g ^ {- 1} x g)
\]

对所有 \(f \in \mathcal{C}(\mathrm{G})\) 及每个 \(x \in \mathrm{G}\) 都成立。它是连续等距表示，而 \(\mathcal{C}(\mathrm{G}^{\sharp})\) 是该表示的不变元素空间。因此，连续投影 \(p = \varrho(1)\) 作用在 \(\mathcal{C}(\mathrm{G})\) 上，且其像为 \(\mathcal{C}(\mathrm{G}^{\sharp})\)（V, 第 466 页推论 3）；其范数满足 \(\leqslant 1\)。

令 \(f \in \mathcal{C}(\mathrm{G}^{\sharp})\) 且 \(\varepsilon > 0\)。存在 \(\widetilde{f} \in \Theta(\mathrm{G})\)，使得 \(\|f - \widetilde{f}\| \leqslant \varepsilon\)（V, 第 462 页命题 4），于是 \(\|f - p(\widetilde{f})\| = \|p(f - \widetilde{f})\| \leqslant \varepsilon\)；由于 \(p(\widetilde{f}) \in \mathcal{C}(\mathrm{G}^{\sharp})\)，我们得到 \(\Theta(\mathrm{G}^{\sharp})\) 在 \(\mathcal{C}(\mathrm{G}^{\sharp})\) 中稠密。

记 R(G) 为 G 在分离的有限维复向量空间中的连续表示所成的格罗滕迪克环（参见 LIE, IX, 第 70 页和 A, VIII, 第 182 页；将其应用于 G 在分离的有限维复向量空间中的连续表示这一加性类，并将其视为 \(\mathbf{C}[G]\)-模）。由于 G 在有限维分离拓扑向量空间中的每个连续线性表示都是半单的（V, 第 457 页命题 1），阿贝尔群 R(G) 是自由的，而不可约酉表示 \(\pi \in \widehat{\mathbf{G}}\) 的类构成一个基（A, VIII, 第 186 页命题 7）。

推论 2. --- a) G 的不可约表示的特征标族是 \(\Theta(G^{\sharp})\) 的一个基；

\begin{enumerate}
\def\labelenumi{\alph{enumi})}
\setcounter{enumi}{1}
\tightlist
\item
  映射 \(u: \mathrm{R}(\mathrm{G}) \otimes_{\mathbf{Z}} \mathbf{C} \to \Theta(\mathrm{G}^{\sharp})\)，满足 \(u(\pi \otimes 1) = \chi_{\pi}\)（对每个 \(\pi \in \widehat{\mathrm{G}}\)），是 \(\mathbf{C}\) 上代数的同构。
\end{enumerate}

第一个断言由命题 9 和推论 1 得出。

由于表示 \(\pi \in \widehat{\mathbf{G}}\) 的类构成自由 \(\mathbf{Z}\)-模 R(G) 的一个基，映射 \(u\) 定义良好。它是 \(\mathbf{C}\)-代数的态射，并由 \(a\) 可知它是同构。

推论 3. --- 令 H 为局部紧拓扑群，且 G 是 H 的紧子群。令 \(\varrho\) 为 H 在希尔伯特空间 E 中的酉表示。若 \(\pi\)-同型分量（在 \(\varrho\) 限制到 G 后）对所有 \(\pi \in \widehat{\mathbf{G}}\) 都是有限维的，则 H 的表示 \(\varrho\) 是半单的，并且 H 的每个不可约酉表示在 \(\varrho\) 中的重数都是有限的。

记 \(\sigma\) 为 \(\varrho\) 限制到 H 的紧子群 G 后的表示。令 \(\pi\) 为 G 的不可约表示。自同态 \(\sigma(\chi_{\pi}) \in \mathcal{L}(\mathrm{E})\) 是有限秩的，因为其像是 \(\overline{\pi}\)-同型分量，即 \(\sigma\) 中的相应子空间（V, 第 465 页推论 2），按假设它是有限维的。因此，自同态 \(\sigma(f)\) 对每个 \(f \in \Theta(\mathrm{G}^{\sharp})\) 都是有限秩的；对每个 \(f \in \mathcal{C}(\mathrm{G}^{\sharp})\)，它都是紧的（V, 第 468 页推论 1 以及 III, 第 4 页命题 2 的推论）。

记 \(j\) 为从 G 到 H 的典范嵌入。这是一个连续映射，并且对于 G 上的每个测度 \(\nu\) 都是 \(\nu\)-真映射（INT, V, 第 69 页，第 6 节第 1 号注 1）。令 \(\mathfrak{B}\) 为 G 中元素 \(e\) 的紧邻域滤子。对于每个 V \(\in\) \(\mathfrak{B}\)，G 上存在一个连续正中心函数 \(f_{\mathrm{V}}\)，其支撑包含于 V 且在 G 上的积分等于 1。令 \(\beta_{\mathrm{V}}\) 为像测度 \(j(f_{\mathrm{V}} \cdot \mu)\)；它是 H 上具有紧支撑的正测度，并满足 \(\varrho(\beta_{\mathrm{V}}) = \sigma(f_{\mathrm{V}})\)（INT, V, 第 71 页，第 6 节第 2 号定理 1）。根据上述结果，在 \(\mathcal{M}_{c}(\mathrm{H})\) 上，由 \(\mathfrak{B}\) 经映射 \(\mathrm{V} \mapsto \beta_{\mathrm{V}}\) 得到的滤子基满足 V, 第 402 页命题 2 的条件，断言得证。

推论 4（外尔等分布判据）

令 M 为 G 上的中心正测度集，满足 \(\nu(G)=1\) 对所有 \(\nu\in M\) 都成立。令 \(\mathfrak{F}\) 为 M 上的滤子。滤子 \(\mathfrak{F}\) 弱收敛到测度 \(\mu^{\sharp}\)，其中测度定义在 \(G^{\sharp}\) 上，其充要条件是：对每个非平凡不可约酉表示 \(\pi\)，都有

\[
\lim _ {\nu , \mathfrak {F}} \int_ {\mathrm{G}} \chi_ {\pi} (x) d \nu (x) = 0.
\]

由于 \(\nu (\mathrm{G}^{\sharp}) = 1 = \mu^{\sharp}(\mathrm{G}^{\sharp})\) 对 \(\nu \in \mathrm{M}\) 成立，上述假设意味着

\[
\lim _ {\nu , \mathfrak {F}} \int_ {\mathrm{G} ^ {\sharp}} \chi_ {\pi} (x) d \nu (x) = \int_ {\mathrm{G} ^ {\sharp}} \chi_ {\pi} (x) d \mu^ {\sharp} (x)
\]

对每个表示 \(\pi\in\widehat{\mathrm{G}}\) 都成立，因而它等价于如下条件

\[
\lim _ {\nu , \mathfrak {F}} \int_ {\mathrm{G} ^ {\sharp}} f (x) d \nu (x) = \int_ {\mathrm{G} ^ {\sharp}} f (x) d \mu^ {\sharp} (x)
\]

由于线性性，对每个函数 \(f \in \Theta(\mathrm{G}^{\sharp})\) 都成立。由于空间 \(\Theta(\mathrm{G}^{\sharp})\) 在 \(\mathcal{C}(\mathrm{G}^{\sharp})\) 中稠密（推论 1），所以该假设等价于滤子 \(\mathfrak{F}\) 收敛到 \(\mu^{\sharp}\)，并且在 \(\mathcal{M}(\mathrm{G}^{\sharp})\) 中收敛；该空间赋予由 \(\mathcal{C}(\mathrm{G})\) 上逐点收敛定义的拓扑。由于 G 紧，这一拓扑与弱收敛拓扑一致（参见 INT, III, 第 59 页，第 1 节第 9 号）。

